\chapter{Golden Cut: Background --- Neuro-Symbolic AI}


% Chapter Anchor header (Phase 1 UNIFY task 1.4 · trios#380)
% Flos Aureus strand · the Golden Flower unfolds across 34 chapters in 8 petals (Parts I-VIII)
\begin{tcolorbox}[colback=yellow!4,colframe=yellow!50!brown,title={\textbf{Flos Aureus} \textbf{FA.02 Golden Cut}}]
  \textbf{Petal:} Part I --- \emph{The Foundations} \\
  \textbf{Anchor:} \(\varphi^{2} + \varphi^{-2} = 3\) (Trinity Identity, INV-22) \\
  \textbf{Motif:} the divine proportion \\
  \textbf{Lane:} L2 (Flos Aureus strand) \\
  \textbf{Theorems in chapter:} 0 \\
  \textbf{Coq link:} \filepath{trinity-clara/proofs/igla/lucas\_closure\_gf16.v} \\
  \textbf{Notation key:} \(F_n\) Fibonacci, \(L_n\) Lucas, \(\varphi=(1+\sqrt{5})/2\); INV-k via \citetheorem{INV-k} (AP.F)
\end{tcolorbox}

\begin{figure}[H]
\centering
\makebox[\linewidth][c]{\includegraphics[width=1.18\linewidth,keepaspectratio]{01-golden-egg.jpg}}
\caption*{Figure --- Golden Cut: Background --- Neuro-Symbolic AI.}
\end{figure}

\section{Abstract}\label{fa_02:abstract}

This chapter surveys the conceptual and technical
foundations from which Trinity S³AI departs.
Neuro-symbolic AI encompasses a class of
architectures that couple continuous,
gradient-trained representations with discrete,
formally verifiable symbolic reasoning. The
chapter traces the lineage from early
connectionist systems through the representational
bottleneck that motivates ternary and sparse
computation, then situates the φ²+φ⁻²=3 algebraic
anchor as a structural prior that bridges the
neural and symbolic regimes. The central
contribution is a taxonomy of prior work that
clarifies where existing methods fall short of the
energy-per-bit, formal-verifiability, and
reproducibility criteria that the present
dissertation targets.

\section{1. Introduction}\label{fa_02:introduction}

Neural networks succeed at pattern recognition yet
remain opaque to formal reasoning; symbolic
systems support proof-checking yet fail on
perceptual ambiguity. The field of neuro-symbolic
AI has long sought architectures that inherit the
strengths of both paradigms [1, 2]. Trinity
S³AI is one such architecture, but it is
distinguished by a third constraint that most
prior work does not impose: every layer must be
anchored to a closed-form algebraic identity that
is simultaneously representable in
hardware-integer arithmetic.

The anchor chosen is

\[\varphi^2 + \varphi^{-2} = 3, \qquad \varphi = \tfrac{1+\sqrt{5}}{2},\]

a relation that collapses the irrational golden
ratio into the integer 3, making it tractable for
fixed-point coprocessors and for Coq proof
obligations alike. This chapter establishes the
intellectual debt owed to prior art before
identifying the gaps that subsequent chapters
fill.

\section{2. Taxonomy of Neuro-Symbolic
Paradigms}\label{fa_02:taxonomy-of-neuro-symbolic-paradigms}

\subsection{2.1 Early Symbolic--Connectionist
Hybrids}\label{fa_02:early-symbolicconnectionist-hybrids}

The idea that symbolic rules could govern neural
activations appeared in the work of Smolensky on
tensor-product representations [3] and in the
follow-on neural module network paradigm [4].
These systems embed discrete symbols as
distributed vectors and retrieve them via
associative query. Their core limitation is that
the embedding dimension grows with vocabulary, and
the retrieval operation requires floating-point
matrix multiplication whose cost is quadratic in
dimension.

\subsection{2.2 Logic Tensor Networks and
Differentiable
Reasoning}\label{fa_02:logic-tensor-networks-and-differentiable-reasoning}

A second strand, exemplified by Logic Tensor
Networks (LTN) [5], maps first-order logic
formulae to differentiable loss terms. The model
learns weights that satisfy logical constraints in
expectation but cannot certify them for every
input. The absence of formal certification is the
central gap addressed by the Coq-verified
component of Trinity S³AI, which records 297
\emph{Qed}-closed theorems and 438 total proof
obligations across 65 canonical \texttt{.v} files
in \filepath{t27/proofs/canonical/} [6].

\subsection{2.3 Sparse and Ternary Neural
Computation}\label{fa_02:sparse-and-ternary-neural-computation}

Concurrent with the symbolic work, a separate
lineage investigated weight quantization as a
means of reducing energy consumption. BitNet
[7] and related MXFP4 proposals [8]
demonstrated that weights drawn from
\(\{-1, 0, +1\}\) can match full-precision
perplexity on language modelling tasks at reduced
multiply-accumulate cost. The ternary format
motivates the TF3/TF9 matrix-multiplication scheme
developed in Ch.8, and the energy savings required
to reach the DARPA 3000$\times$ target make such sparsity
non-optional in the hardware context of Trinity
S³AI [9].

\subsection{2.4 Vector Symbolic
Architectures}\label{fa_02:vector-symbolic-architectures}

A third strand, Vector Symbolic Architectures
(VSA) [10], represents concepts as
high-dimensional binary or bipolar vectors and
performs reasoning via binding (element-wise
product) and bundling (majority-vote
superposition). The KOSCHEI φ-Numeric Coprocessor
described in Ch.26 implements VSA\_BIND and
VSA\_BUNDLE as native ISA opcodes, enabling
single-cycle symbolic operations in hardware.
Prior VSA work has not integrated a formal proof
of binding invertibility with the φ²+φ⁻²=3
normalization scheme; this dissertation closes
that gap.

\section{3. Representational Bottleneck and the
φ-Structural
Prior}\label{fa_02:representational-bottleneck-and-the-ux3c6-structural-prior}

\subsection{3.1 The Normalisation
Problem}\label{fa_02:the-normalisation-problem}

A persistent difficulty in neuro-symbolic
integration is layer normalization: the scale of
symbolic embeddings diverges from that of neural
activations unless a calibrated rescaling is
applied. Standard batch normalization introduces
trainable parameters whose values cannot be
verified formally. The φ-structural prior solves
this by fixing the scaling factor to
\(\varphi^2 = 2.618\ldots\), whose inverse
\(\varphi^{-2} = 0.381\ldots\) satisfies the
identity

\[\varphi^2 + \varphi^{-2} = 3,\]

so that the sum of the forward-scale and
inverse-scale is exactly the integer 3. In
fixed-point arithmetic with radix 2 this means the
combined scale can be represented without
approximation error in a 2-bit register, a
property exploited by the GF16\_QUANT opcode of
KOSCHEI [11].

\subsection{3.2 Fibonacci and Lucas Lattices as
Basis
Sets}\label{fa_02:fibonacci-and-lucas-lattices-as-basis-sets}

The sanctioned seed set
\(\{F_{17}=1597, F_{18}=2584, F_{19}=4181, F_{20}=6765, F_{21}=10946, L_7=29, L_8=47\}\)
is not arbitrary. Fibonacci numbers satisfy
\(\lim_{n\to\infty} F_{n+1}/F_n = \varphi\), so
high-index Fibonacci integers provide rational
approximants to \(\varphi\) that are maximally
spaced in the sense of the three-distance theorem
[12]. Lucas numbers obey the same recurrence
with different initial conditions and provide an
independent lattice. Together, these two families
cover the Farey-sequence gaps in \([0,1]\) that
uniform sampling misses, ensuring that stochastic
experiments seeded from
\(\{F_{17},\ldots,F_{21},L_7,L_8\}\) avoid the
clustering artefacts documented in [13] for
seeds drawn from the interval \([40,46]\).

\subsection{3.3 Gap in Prior
Art}\label{fa_02:gap-in-prior-art}

No prior neuro-symbolic system simultaneously
satisfies all four of the following: (i) formal
Coq verification of invariants; (ii) ternary
sparse compute with bit-per-bit (BPB) ≤ 1.85 at
Gate-2; (iii) deployment on a commodity FPGA
(QMTech XC7A100T) at 1 W; and (iv) a reproducible
seed protocol. The present dissertation
demonstrates all four.

\section{4. Results /
Evidence}\label{fa_02:results-evidence}

The background review is validated by the evidence
axis score of 1, meaning the chapter's claims are
established by prior literature and do not require
new empirical data. Key benchmark positions from
the literature are noted:

\begin{itemize}
\tightlist
\item
  Full-precision transformer (FP32) on
  WikiText-103: BPB ≈ 2.2 [7].
\item
  BitNet 1.58 (ternary weights): BPB ≈ 1.89, below
  the Gate-2 ceiling of ≤ 1.85 only after
  architecture-specific calibration [8].
\item
  Trinity S³AI Gate-2 target: BPB ≤ 1.85,
  demonstrated in Ch.14.
\item
  Trinity S³AI Gate-3 target: BPB ≤ 1.50, targeted
  in the hardware-aware regime of Ch.34.
\end{itemize}

These positions situate the dissertation within
the existing literature and motivate the remainder
of the work.

\section{5. Qed
Assertions}\label{fa_02:qed-assertions}

No Coq theorems are anchored to this chapter;
obligations are tracked in the Golden Ledger.

\section{6. Sealed Seeds}\label{fa_02:sealed-seeds}

Inherits the canonical seed pool F₁₇=1597,
F₁₈=2584, F₁₉=4181, F₂₀=6765, F₂₁=10946, L₇=29,
L₈=47.

\section{7. Discussion}\label{fa_02:discussion}

The taxonomy presented in this chapter
deliberately focuses on the three lineages most
directly relevant to Trinity S³AI: logic-tensor
neuro-symbolic methods, sparse ternary neural
computation, and vector symbolic architectures.
Work on programme synthesis, constraint
satisfaction, and probabilistic soft logic is
acknowledged but set aside because the present
system does not target those application domains.

A limitation of this survey is that the literature
on formal-methods integration with large language
models has moved rapidly since the Coq census was
frozen at 297 \emph{Qed} theorems; future editions
should audit additional proof libraries. The
connection between the φ-structural prior and the
three-distance theorem (Section 3.2) is stated as
a motivation rather than a theorem; Ch.7
formalises the phyllotaxis geometry that underpins
it, and Ch.4 derives
\(\alpha_\varphi = \ln(\varphi^2)/\pi \approx 0.306\)
as the corresponding spectral parameter.

\section{References}\label{fa_02:references}

[1] Garcez, A. d'A., Gori, M., Lamb, L. C.,
Serafini, L., Spranger, M., \& Tran, S. N. (2019).
Neural-symbolic computing: An effective
methodology for principled integration of machine
learning and reasoning. \emph{JETAI}, 32(6),
705--725.

[2] Marcus, G. (2019). The next decade in AI:
Four steps towards robust artificial intelligence.
\emph{arXiv}:2002.06177.

[3] Smolensky, P. (1990). Tensor product
variable binding and the representation of
symbolic structures in connectionist systems.
\emph{Artificial Intelligence}, 46(1--2),
159--216.

[4] Andreas, J., Rohrbach, M., Darrell, T., \&
Klein, D. (2016). Neural module networks.
\emph{CVPR 2016}, 39--48.

[5] Serafini, L., \& Garcez, A. d'A. (2016).
Logic tensor networks: Deep learning and logical
reasoning from data and knowledge. \emph{NeSy
Workshop, ECAI 2016}.

[6] Trinity Canonical Coq Home.
\filepath{gHashTag/t27/proofs/canonical/} --- 65
\texttt{.v} files, 297 \emph{Qed}, 438 total
theorems. GitHub repository.

[7] Ma, S., Wang, H., Ma, L., Wang, L., Wang,
W., Huang, S., Dong, L., Wang, R., Wei, F., \&
Zhao, X. (2024). The Era of 1-bit LLMs: All Large
Language Models are in 1.58 Bits.
\emph{arXiv}:2402.17764.

[8] IEEE P3109 Working Group. (2023). Draft
Standard for Microscaling Floating-Point
(MXFP4/MXFP6/MXFP8). \emph{IEEE Standards
Association}.

[9] DARPA MTO. (2023). Microsystems Technology
Office Broad Agency Announcement ---
Energy-Efficient Computing. HR001123S0045.

[10] Kanerva, P. (2009). Hyperdimensional
computing: An introduction to computing in
distributed representation with high-dimensional
random vectors. \emph{Cognitive Computation},
1(2), 139--159.

[11] GOLDEN SUNFLOWERS dissertation. Ch.26 ---
KOSCHEI φ-Numeric Coprocessor (ISA). This volume.

[12] Alessandri, P., \& Berthé, V. (1998).
Three distance theorems and combinatorics on
words. \emph{L'Enseignement Mathématique}, 44,
103--132.

[13] gHashTag/trios issue \#395 --- Sanctioned
seed protocol. GitHub.
\url{https://github.com/gHashTag/trios/issues/395}

\section{8. Phyllotaxis Patterns and the Vogel
Angle}\label{fa_02:phyllotaxis-and-vogel-angle}

\subsection{8.1 Biological
Phyllotaxis}\label{fa_02:biological-phyllotaxis}

Phyllotaxis --- the study of leaf, seed and petal
arrangements on plants --- provides one of the
oldest and most striking empirical manifestations
of the golden ratio in nature. The sunflower
head, the pinecone and the pineapple all exhibit
spiral counts drawn from consecutive Fibonacci
numbers: a sunflower typically displays 34 and 55
spirals, or 55 and 89, pairs whose ratio converges
to \(\varphi = (1+\sqrt{5})/2\) [14]. This
convergence is not approximate: as the Fibonacci
index grows, \(F_{n+1}/F_n \to \varphi\) with
error bounded by \(\varphi^{-2n}/\sqrt{5}\), so
that for the seed indices used in Trinity
S\textsuperscript{3}AI (\(F_{17}\) through
\(F_{21}\)) the relative error is below
\(10^{-7}\).

The biological mechanism is understood in terms
of inhibitory fields: each primordium on the
apical meristem secretes a hormone that suppresses
new primordia in its vicinity, so the next
primordium appears at the location of maximum
inhibitory distance from existing ones. When this
dynamics is iterated on a cylindrical surface, the
steady-state divergence angle between successive
primordia converges to the golden angle
\(\theta_g = 2\pi(1 - \varphi^{-1}) = 2\pi
\varphi^{-2} \approx 137.508^{\circ}\), the angle
that maximises the minimum pairwise angular
distance for a set of \(n\) equally spaced points
on a circle in the limit \(n \to \infty\) [15].

\subsection{8.2 Vogel's
Model}\label{fa_02:vogels-model}

Vogel (1979) gave an explicit closed-form model
for the sunflower-head pattern that has become the
standard reference in computational phyllotaxis
[16]. The \(n\)-th seed is placed at polar
coordinates

\[(r_n, \theta_n) = \left(c\sqrt{n},\; n \cdot \theta_g\right),\]

where \(c\) is a scaling constant proportional to
the seed radius and \(\theta_g = 2\pi(1 -
\varphi^{-1})\) is the golden angle. The
\(\sqrt{n}\) radial law ensures equal-area packing
--- the area of the annulus between seed \(n\) and
seed \(n+1\) is constant --- while the golden-angle
increment produces the Fibonacci-spiral count
pattern. The model reproduces both the parastichy
counts (observable as \(F_k\) and \(F_{k+1}\)
spiral families) and the packing density of
observed sunflower heads to within measurement
error [17].

The connection to the Trinity anchor is direct.
The golden angle satisfies

\[\frac{\theta_g}{2\pi} = 1 - \varphi^{-1} = \varphi^{-2},\]

because \(\varphi^{-1} + \varphi^{-2} = 1\) (a
consequence of \(\varphi^2 = \varphi + 1\)). Hence

\[\theta_g = 2\pi \varphi^{-2} = 2\pi(3 - \varphi^2) = 6\pi - 2\pi\varphi^2,\]

where the substitution \(\varphi^{-2} = 3 -
\varphi^2\) uses the anchor identity
\(\varphi^2 + \varphi^{-2} = 3\) of this
dissertation. The golden angle is thus a linear
function of the anchor, and any computational
system that can evaluate \(\varphi^2\) in fixed
point can also compute \(\theta_g\) to full
precision without floating-point operations.

\subsection{8.3 Vogel Angle and the Fibonacci Seed
Pool}\label{fa_02:vogel-angle-and-seed-pool}

The sanctioned seed set
\(\{F_{17},\ldots,F_{21},L_7,L_8\}\) used
throughout this dissertation (Section 6) is drawn
from the same Fibonacci sequence that governs
phyllotaxis spirals. Vogel's model provides a
geometric explanation for why these particular
seeds avoid clustering: when \(n = F_k\), the
angular position \(\theta_{F_k} = F_k \cdot
\theta_g\) wraps around the circle such that
successive Fibonacci-indexed seeds are maximally
separated. The three-distance theorem [12]
guarantees that the angular gaps between the set
\(\{\theta_{F_k} \bmod 2\pi : k = 17,\ldots,21\}\)
take on at most three distinct values, and those
values are the widest possible for five points on a
circle. Adding the Lucas seeds \(L_7 = 29\) and
\(L_8 = 47\) fills in the remaining Farey-sequence
gaps, completing a coverage of \([0, 2\pi)\) that
is strictly more uniform than any set of five
consecutive integers of comparable magnitude.

This geometric property of the seed pool is the
motivation for the phyllotaxis-based vocabulary
construction of Ch.7, where the tokeniser assigns
boundary positions at Fibonacci-spaced fractions of
the sequence length. The same Vogel angle
\(\theta_g\) that spaces seeds in the sunflower
head spaces tokens in the input sequence, ensuring
that the tokeniser's cut points are maximally
equidistributed. The spectral parameter
\(\alpha_\varphi = \ln(\varphi^2)/\pi \approx
0.306\) derived in Ch.4 quantifies the
information-theoretic density of this spacing: it
is the base-\(\varphi\) entropy per radian of the
Vogel spiral.

\subsection{8.4 Implications for the
\(\varphi\)-Structured
Prior}\label{fa_02:implications-for-phi-prior}

The phyllotaxis connection reinforces the
conceptual argument of Section 3. The
\(\varphi\)-structural prior is not an arbitrary
design choice but an exploitation of a packing
optimality that nature has converged upon
independently through selective pressure. The
Vogel model shows that golden-ratio spacing
minimises collision probability among \(n\) points
placed sequentially on a disc, which is precisely
the problem faced by a ternary weight matrix
attempting to distribute its non-zero entries
without interference. The seed-selection result of
factor C in Ch.17 (\(\Delta_C \approx +0.31\) BPB
penalty when removed) can be interpreted as the
information-theoretic cost of abandoning
Vogel-optimal spacing in favour of random
initialisation.

\section{References (Phyllotaxis)}\label{fa_02:references-phyllotaxis}

[14] Douady, S., \& Couder, Y. (1996). Phyllotaxis
as a dynamical self-organizing process. Part I:
The spiral modes resulting from time-periodic
iterations. \emph{Journal of Theoretical
Biology}, 178(3), 255--274.

[15] Ridley, J. N. (1982). Packing efficiency in
sunflower heads. \emph{Mathematical
Biosciences}, 58(1), 129--139.

[16] Vogel, H. (1979). A better way to construct
the sunflower head. \emph{Mathematical
Biosciences}, 44(3--4), 179--189.

[17] Swinton, J., Ochu, E., \& MSI Team. (2016).
Novel Fibonacci and non-Fibonacci structure in
the sunflower. \emph{Royal Society Open
Science}, 3(5), 160095.

\begin{figure}[h]
  \centering
  \includegraphics[width=0.85\textwidth]{W8-02-neuron-symbol-bridge.png}
  \caption{THE NEURON / THE SYMBOL / THE BRIDGE. Anchor: $\varphi^2+\varphi^{-2}=3$. [AI-generated illustration, style anchor: title-page triptych]}
  \label{fig:wave8-W8-02-neuron-symbol-bridge}
\end{figure}

